\section{Tele-Educação e Capacitação Continuada Descentralizada}
\label{sec:tele-educacao}

A topologia em rede dos gêmeos digitais transborda as fronteiras da graduação, redefinindo os arcabouços de educação continuada e treinamento especializado para cirurgiões-dentistas formados. Especialmente em cenários geoeconômicos marcados por extensões continentais e assimetrias na distribuição de especialistas, como ocorre na malha demográfica do Brasil, a tele-educação adquire o status de infraestrutura de saúde pública~\cite{frota2025sus}.

O conceito de \textit{Tele-mentoring} háptico inaugura a era do treinamento cirúrgico supervisionado à distância. Por meio de conexões 5G de ultrabaixa latência (uRLLC) combinadas a gêmeos digitais sincronizados bilateralmente, um endodontista experiente situado em um centro de excelência metropolitano pode "sentir" a mesma resistência tátil de um operador generalista que manipula um paciente virtual ou atende um paciente remoto em uma localidade vulnerável, guiando milimetricamente suas mãos ou corrigindo vetores de inclinação em tempo real~\cite{cuadros2023access}.

\destaque{A capacitação continuada baseada em simulação distribuída é uma alavanca para a equidade sistêmica, atenuando a obsolescência profissional precoce e capilarizando a Odontologia Baseada em Evidências para os rincões do sistema de saúde.}

Para profissionais inseridos em ecossistemas marginalizados ou programas públicos (como o SUS no Brasil), a simulação remota contorna restrições de tempo, viagens e perda de dias úteis, propiciando programas de \textit{lifelong learning} gamificados e assíncronos. Essa arquitetura tem provado ser uma estratégia robusta não apenas na transferência de competências cognitivas de diagnóstico, mas em protocolos cirúrgicos emergenciais, como traqueostomias, manejo de avulsões ou controle de hemorragias arteriais maxilares complexas, aumentando radicalmente a autoconfiança e as taxas de sucesso clínico do profissional periférico~\cite{du2020marginalized}.
